\section*{Conventions and manuscript concordance}
\begingroup\fontsize{9.3}{10.8}\selectfont
\setlength{\tabcolsep}{4pt}
\begin{longtable}{p{.19\textwidth}p{.48\textwidth}p{.27\textwidth}}
\toprule Our notation & Manuscript notation and meaning & Passage\\\midrule
$E,P,Q;T$ & $\Pi,\Pi_0,G$ in Part A; finite extension, surface subgroup, quotient, and orientation line & \S52, ff.737--743, scans179--191\\
$\mathcal L(E),LT,TL$ & Virtual loop structures; the defining commensurability class is retained for LT/TL & \S52, scans180--191\\

$\Pi_0\subset\Pi$ & Discrete group and its specified profinite completion & \S52, f.750, scan205\\
$\mathcal G_0,\mathcal G$ & Discrete geometric extension and its closure; the closure need not be its full profinite completion & \S52, ff.750--751, scans205--207\\
$\mathcal M,\Gamma$ & Ambient profinite group and $\mathcal M/\mathcal G$ & \S52, ff.750--751, scans205--207\\
$\chi$ & Common cyclotomic character with open image & \S52, f.751, scan207\\
$\mathbb A_f$ & $\mathcal A_\Q=\widehat\Z\otimes\Q$ & \S52, f.748, scan201\\
$G,\rho$ & $\underline{\mathcal M}_\Q$, a finite-type algebraic target of an admissible adelic representation & \S52, ff.751--752, scans208--210\\
$E_\Gamma,\Gamma_{nM}$ & Schematic envelope and its virtually abelian finite stages; $M$ is a normal lattice & \S50, scans25--31\\
$V(F),U_n,L,\gamma_n$ & Virtually abelian system, class-$n$ unipotent quotient, graded Lie algebra, lower central series; auxiliary constructions & Modern notation; B.2--B.9\\
$B,S_3,\varphi,\iota$ & $C_2*C_3$, its finite quotient, the quotient map and the distinct level-two inclusion & \S47, ff.415--420v; cote140-3 scans229--240\\
$\widetilde{\SL}_2(\Z)$ & Inverse image of $\SL_2(\Z)$ in the universal covering group of $\SL_2(\R)$ & \S53, ff.762--764, scans229--233\\
$z=(\sigma_1\sigma_2)^3$ & Central braid; the kernel over $\SL_2(\Z)$ is generated by $z^2$ & \S53, ff.764--766, scans233--237\\
$ST^+_{1,1},SN_{1,1}$ & Positive special Teichm\"uller group and arithmetic extension with specified tangent & \S53, scans222--241\\
$\operatorname{Cov}(S,\Phi)$ & Category of coverings of a germ; not a tower of Lie completions & \S53.5, ff.768--769, scans242--243\\
$H^1,H^1_f,U_n/F^0$ & Section classes, Selmer classes, and local de Rham quotient; local conditions are part of the data & Modern notation; Part F\\
$K,E,P,\Delta,G_1,s$ & $\Q(i)$, $y^2=x^3-x$, its pro-$\ell$ group, abelian-field kernel, torsion-field kernel, tangent section & Modern CM notation; Part G\\
\bottomrule
\end{longtable}
\endgroup
\smallskip
Unless specified otherwise, groups are discrete; profinite homomorphisms
are continuous. In characteristic zero, unipotent completion means the
universal prounipotent affine group. A genuine Galois action includes a
section or tangential base point. We use $\gamma_1F=F$ and
$\gamma_{r+1}F=[F,\gamma_rF]$; degree-$m$ derivations send $L_1$ to
$L_{m+1}$, independently of Tate-weight conventions. Locally bound symbols
are defined at first use. Repeated letters are not identified across Parts
without an explicit comparison.

Transcription marks, including doubtful and deleted readings, are evidence
to inspect, not axioms. The readings ledger gives mathematical evidence
without claiming a final diplomatic reading of the handwriting.
\clearpage
